\subsection{System}
For computing the usable space, deduct the overhead \textit{after} multiplying with the free fraction,
e.g. for a block size of 4kB (with kB actually being KiB), header and directories having a combined size of 100 bytes and \texttt{PCTFREE} set to 20\%,
the available space is $\ceil{4096 \cdot 0.8 - 100} = 3177$.

To compute the maximum number of tuples in a slotted page, remember that it also has a header with a certain number of bytes for the slotting.
Add that to the number of bytes per tuple when computing, even though it doesn't count to the tuple size!

Row stores are better if we need to access many columns, column stores otherwise. This is (primarily) due to indexing.


\subsubsection{Indexing}
The global depth is the maximum Local Depth of all buckets.

The hash key is the actual value of the column to be hashed typically.
Thus, an assignment may be very misleading in that the block ID and hash key are asked.
The hash key is the value before hashing, the block ID afterwards is the block the hashing algorithm assigned it to.

When splitting a block, its $N$ entries must be repartitioned. There are $2^N$ possible ways in which to do this.
So, there is a $\frac{2}{2^N}$ probability for there to be a \bi{recursive overflow}, an overflow in which all keys end up in the same bucket again
and need further splitting.
